\section{Experimental Setup}
\label{sec:experiments}

\subsection{SENTINEL-DB}

SENTINEL-DB is a $\sim$85\,GB foundational dataset harmonizing water quality
data from thirteen public sources totaling over 390 million real records plus
291K sensor sequences, 2,986 satellite tiles, 268K ecotoxicity records, and
5,000 behavioral trajectories. The largest contributor is NEON Aquatic
(351.7M rows of continuous high-frequency sensor data from 34 aquatic sites,
4 WQ products, 24 months). The discrete WQ data comprises 17.99M harmonized
river quality records from GRQA v1.3 (22 parameters, 105 countries, 94K
sites), 18.27M discrete samples from the EPA Water Quality Portal (18 HUC2
basins), and 291K sensor time series from 1,130 USGS NWIS stations.
Additionally, 268,029 dose-response records from EPA ECOTOX (1,391
chemicals, 8 contaminant classes) support molecular contaminant
classification. A unified parameter ontology and H3 hexagonal spatial
indexing (resolution 8) enable cross-source deduplication and satellite
co-registration.

\subsection{Per-Encoder Experiments}

\paragraph{AquaSSM.}
Pretrained via Masked Parameter Prediction (MPP) on 20,000 real USGS NWIS
instantaneous-value sequences (5 parameters: DO, pH, specific conductance,
temperature, turbidity) from 1,115 monitoring stations (291K sequences
available across 1,130 stations). Fine-tuned for anomaly detection using
EPA threshold-based labels (e.g., DO $<$ 4 mg/L, pH outside 6.5--9.0).
Sequence length $T=128$, batch size 32, learning rate $5 \times 10^{-4}$
with cosine decay.

\paragraph{HydroViT.}
MAE-pretrained on 2,986 real Sentinel-2 L2A tiles (10 spectral bands,
224$\times$224 pixels) from Planetary Computer. Water quality regression
head fine-tuned on 4,202 co-registered satellite--in situ WQ pairs from
GRQA (2015--2020) and EPA WQP spanning 170 geographic cells across 105
countries, plus USGS NWIS stations with GPS-precise ($\sim$10\,m)
coordinates and tightly co-registered Sentinel-2 tiles ($\pm$1.5\,km,
$\pm$3\,days). Split: 2,941 train / 630 val / 631 test. Training:
masking ratio 75\%, head LR $3 \times 10^{-4}$, backbone LR $5 \times
10^{-5}$, 200 epochs total.

\paragraph{MicroBiomeNet.}
Trained on 20,288 real 16S rRNA OTU samples from the Earth Microbiome
Project (EMP), covering 8 aquatic environment classes per the EMPO\_3
ontology. CLR-transformed compositional features, Aitchison-geometry
attention layers, gradient clamping ($\pm 5$). Learning rate $3 \times
10^{-4}$ with cosine decay over 100 epochs.

\paragraph{ToxiGene.}
Trained on 4 real transcriptomic datasets from NCBI GEO (zebrafish PAH,
copper, Daphnia, fathead minnow; $\sim$84K genes total) augmented with
268K dose-response records from EPA ECOTOX (8 contaminant classes: heavy
metals, pesticides, pharmaceuticals, PAHs, PCBs, PFAS, nanomaterials,
nutrients). P-NET biological hierarchy with sparse Reactome adjacency
matrices. 80 epochs with cosine learning rate schedule.

\paragraph{BioMotion.}
Phase 1: diffusion denoising pretraining on 17,074 real Daphnia behavioral
toxicity tests from EPA ECOTOX (noise schedule $\beta \in [10^{-4},
0.02]$). Each test encodes a concentration-response curve across behavioral
endpoints (locomotion, swimming velocity, equilibrium, activity, motility)
as a 200-timestep trajectory. Phase 2: anomaly classification fine-tuning
distinguishing toxicant-exposed responses (LOEC/EC50-level effects:
locomotion impairment, loss of equilibrium, reduced swimming velocity) from
sub-threshold controls. Normal/anomaly labels derived from ECOTOX endpoint
codes (NOEC/EC0 $\rightarrow$ normal; LOEC/EC50/EC100 $\rightarrow$ anomaly).

\subsection{Fusion Evaluation}

\paragraph{31-Condition Ablation.}
All $2^5 - 1 = 31$ non-empty subsets of five modalities evaluated on
held-out pollution events.

\paragraph{Missing Modality Robustness.}
Performance measured when modalities are randomly dropped at test time,
producing degradation curves.

\paragraph{Detection Lead Time.}
For events where all modalities are available, we compare how far in
advance each individual modality versus the fused system detects the event.

\paragraph{ToxiGene Downstream Validation.}
Four additional experiments assess ToxiGene's operational reliability:
(1)~\emph{False positive rate}: inference on 10 clean NEON reference sites
to measure FPR under baseline conditions;
(2)~\emph{Parameter attribution}: gradient-based attribution through the
P-NET hierarchy to recover known gene--pathway--contaminant associations;
(3)~\emph{Temporal persistence}: signal-to-noise ratio between contaminated
and clean temporal windows on held-out time series;
(4)~\emph{Pollution fingerprinting}: discriminating contaminant mixtures
from single-contaminant exposures on held-out samples.

\subsection{Implementation Details}

All encoders project to a shared 256-dimensional embedding space.
Training uses AdamW with cosine learning rate decay. Models are
implemented in PyTorch 2.0 and trained on NVIDIA A100 80GB GPUs.
Total parameter count: 189.5M across all encoders plus 33.6M in the
Perceiver IO fusion layer. Code and trained checkpoints are available
at \url{https://github.com/bcheng/sentinel}.
